\documentclass[11pt]{article}

\usepackage[margin=1in]{geometry}
\usepackage{booktabs}
\usepackage{amsmath}
\usepackage{amssymb}
\usepackage{array}
\usepackage{longtable}
\usepackage{hyperref}
\usepackage{xcolor}
\usepackage{enumitem}

\hypersetup{
    colorlinks=true,
    linkcolor=blue,
    urlcolor=blue,
    citecolor=blue
}

\setlist[itemize]{leftmargin=1.2em}

\title{Stage 1 Self-Test Report: Model Attempts and Held-Out Performance}
\author{Aaron Gao}
\date{\today}

\begin{document}
\maketitle

\section{Purpose of the Self-Test}

This report documents the self-test methodology for all major Stage 1 model
attempts. The grading requirement is that the provided data must be split into
training, validation, and test sets, and that model performance must be reported
on the held-out test set. The grading also emphasizes two points:
\begin{itemize}
    \item the methodology must be sound: proper chronological splitting, no
    leakage, and sensible validation;
    \item the reported test performance should exceed the provided baseline.
\end{itemize}

The final submitted Stage 1 model is
\texttt{mlp/day3\_stage1\_recent\_window}. However, this report also explains
the self-test evidence for the other models I tried: the 3-day XGBoost
baseline, enhanced Ridge, Stage 1 XGBoost, base MLP, tuned MLP, LSTM,
style-dynamic MLP, recent-window ensemble, optimized portfolio branch,
model-switch router, and excess-target MLP branch.

\section{Prediction Target and Test Metric}

Stage 1 uses a 3-trading-day holding horizon. For stock \(i\) on trading date
\(t\), the supervised target is
\[
    y_{i,t}
    =
    r_{i,t}^{(3)}
    =
    \frac{P_{i,t+3}}{P_{i,t}} - 1,
\]
where \(P_{i,t}\) is the adjusted close price. The benchmark is the CSI500
index return over the same horizon:
\[
    R_{b,t}^{(3)}
    =
    \frac{I_{t+3}}{I_t} - 1.
\]
For a portfolio with weights \(w_{i,t}\), the realized portfolio return is
\[
    R_{p,t}^{(3)}
    =
    \sum_i w_{i,t} r_{i,t}^{(3)}.
\]
The primary held-out test metric is mean 3-day excess return:
\[
    \overline{R}_{excess}
    =
    \frac{1}{T}\sum_{t=1}^{T}
    \left(R_{p,t}^{(3)} - R_{b,t}^{(3)}\right).
\]
I also monitor positive excess rate and rank information coefficient (rank IC).
Positive excess rate is the fraction of test dates on which the portfolio beats
the CSI500 benchmark. Rank IC is the average daily Spearman correlation between
predicted stock scores and realized 3-day returns.

\section{Common Split and Leakage-Control Rules}

All model comparisons use chronological splits by trading date. I do not split
randomly by rows, because stocks on the same date share market conditions and
random row splitting would mix future and past regimes. All stocks from the
same date are assigned to the same split.

\subsection{Canonical Split}

The main canonical self-test split is shown in Table
\ref{tab:canonical-split}. This is the split used for the main comparison table
across model families.

\begin{table}[h]
\centering
\caption{Canonical chronological train / validation / test split.}
\label{tab:canonical-split}
\begin{tabular}{lrrl}
\toprule
Split & Rows & Date range & Usage \\
\midrule
Training & 109,870 & up to 2026-03-09 & Fit model parameters \\
Validation & 4,961 & 2026-03-17 to 2026-03-30 & Select hyperparameters and portfolio rule \\
Test & 4,919 & 2026-04-08 to 2026-04-21 & Final held-out self-test \\
\bottomrule
\end{tabular}
\end{table}

The 5-trading-day embargo between train, validation, and test blocks reduces
overlap from the 3-day forward-return labels. The test set is not used to pick
model architecture, training window, tree depth, regularization, top-\(k\), or
weighting rule.

\subsection{Robust Split}

For the strongest candidates, I also used a longer robust split. This checks
whether a model survives a longer held-out period rather than only one favorable
10-day test block.

\begin{table}[h]
\centering
\caption{Robust chronological train / validation / test split.}
\label{tab:robust-split}
\begin{tabular}{lrrl}
\toprule
Split & Rows & Date range & Usage \\
\midrule
Training & 102,456 & up to 2026-02-06 & Fit model parameters \\
Validation & 7,402 & 2026-02-24 to 2026-03-16 & Select hyperparameters and portfolio rule \\
Test & 9,871 & 2026-03-24 to 2026-04-21 & Robust held-out self-test \\
\bottomrule
\end{tabular}
\end{table}

\subsection{Leakage Controls}

The following leakage controls are used throughout the experiments:
\begin{itemize}
    \item Splits are chronological by date, never random by row.
    \item All stocks from the same trading date stay in the same split.
    \item The 3-day target uses \(P_{i,t+3}\), so rows without observable
    future returns are removed from supervised training.
    \item A 5-trading-day embargo separates training, validation, and test
    blocks.
    \item Hyperparameters are selected on validation only.
    \item The test set is used only for final reporting.
    \item Portfolio weights are generated from model scores and then normalized
    under the competition constraints: long-only, total weight 1, maximum
    single-name weight 10\%, and at least 30 positive-weight stocks.
\end{itemize}

\section{Validation Procedure by Model Attempt}

Table \ref{tab:validation-procedure} summarizes how each attempted model used
the training and validation data before being evaluated on the held-out test
set.

{\small
\setlength{\tabcolsep}{4pt}
\begin{longtable}{p{0.19\textwidth}p{0.36\textwidth}p{0.35\textwidth}}
\caption{Validation procedure for each major Stage 1 model attempt.}
\label{tab:validation-procedure}\\
\toprule
Model attempt & What was trained on the training set & What was selected on validation \\
\midrule
\endfirsthead
\toprule
Model attempt & What was trained on the training set & What was selected on validation \\
\midrule
\endhead
3-day XGBoost baseline &
Gradient-boosted tree model trained on public price-volume technical factors
to predict raw 3-day forward return. &
Validation selected a tradable top-ranked portfolio and established the
provided baseline for comparison. \\
\midrule
Enhanced Ridge &
Ridge regression trained on an expanded 3-day factor panel including
market-relative returns, intraday range, volatility, liquidity, beta, and
cross-sectional rank features. &
Validation selected the Ridge regularization strength and score-to-weight
portfolio rule. \\
\midrule
Stage 1 XGBoost &
XGBoost trained on the enhanced public-price factor panel. &
Validation selected tree regularization profile, recency weighting, top-\(k\),
and weight transformation. \\
\midrule
Base MLP &
MLP regressor trained on standardized enhanced factors with ReLU activation,
Adam optimization, early stopping, and \(L_2\) regularization. &
Validation selected top-\(k\) and capped score/softmax weighting. \\
\midrule
Tuned MLP &
Same model family as base MLP, with a narrower architecture and regularization
search around the best MLP settings. &
Validation selected architecture, softmax temperature, risk-aware softmax
variant, and top-\(k\). \\
\midrule
LSTM &
Sequence model trained on recent 20-day feature histories for each stock. &
Validation selected the sequence model configuration and the final portfolio
mapping. \\
\midrule
Style-dynamic MLP &
MLP with market-state-aware portfolio behavior using trend, volatility,
drawdown, breadth, and style/regime signals. &
Validation selected dynamic allocation rules, including top-\(k\) and weighting
changes across regimes. \\
\midrule
Recent-window MLP &
MLP trained on recent rolling windows rather than always using full history. &
Validation selected lookback \(\{63,126,189,252,\text{full}\}\), architecture,
top-\(k\), and capped softmax/risk-softmax weighting. \\
\midrule
Recent-window ensemble &
Combination of recent-window variants trained on the same leakage-controlled
split. &
Validation selected how to combine recent-window scores and how to map ensemble
scores into weights. \\
\midrule
Optimized portfolio branch &
Prediction model similar to the MLP branch, with more emphasis on the portfolio
construction layer. &
Validation selected score-to-weight transformations and concentration/risk
controls. \\
\midrule
Model-switch router &
Separate recent-window and style-dynamic experts were trained. A router used
market-state features to choose between experts. &
Validation and walk-forward windows selected the router policy. The objective
was average excess return with penalties for volatility and negative windows. \\
\midrule
Excess-target MLP &
Recent-window MLP trained with labels transformed from raw 3-day return to
3-day excess return, with winsorization or z-score target denoising. &
Validation selected target transform, lookback, architecture, top-\(k\), and
weighting method. \\
\bottomrule
\end{longtable}
}

\section{Canonical Held-Out Test Results}

Table \ref{tab:canonical-results} reports the canonical held-out test result
for every major model attempt with an available canonical self-test result. The
reported metric is mean 3-day excess return after constructing the final
portfolio. The provided 3-day XGBoost baseline is the reference point.

\begin{table}[h]
\centering
\caption{Canonical held-out test performance by model attempt.}
\label{tab:canonical-results}
\begin{tabular}{p{0.34\textwidth}rrp{0.25\textwidth}}
\toprule
Model attempt & Test excess & Gain vs. baseline & Test-set conclusion \\
\midrule
Provided 3-day XGBoost baseline & \(+0.767\%\) & -- & Baseline to beat \\
Enhanced Ridge & \(+1.273\%\) & \(+0.506\%\) & Exceeds baseline \\
Stage 1 XGBoost & \(+1.394\%\) & \(+0.627\%\) & Exceeds baseline \\
Base Stage 1 MLP & \(+1.741\%\) & \(+0.974\%\) & Exceeds baseline \\
Tuned MLP & \(+1.744\%\) & \(+0.977\%\) & Exceeds baseline \\
LSTM & \(+0.890\%\) & \(+0.123\%\) & Exceeds baseline, but weak \\
Style-dynamic MLP & \(+0.308\%\) & \(-0.459\%\) & Does not exceed baseline \\
Recent-window ensemble & \(+1.244\%\) & \(+0.477\%\) & Exceeds baseline \\
Optimized portfolio branch & \(+1.424\%\) & \(+0.657\%\) & Exceeds baseline \\
Model-switch router & \(+1.842\%\) & \(+1.075\%\) & Exceeds baseline \\
Excess-target weighted MLP & \(+0.993\%\) & \(+0.226\%\) & Exceeds baseline, but not robust \\
\textbf{Recent-window MLP} & \(\mathbf{+2.526\%}\) & \(\mathbf{+1.759\%}\) & \textbf{Best canonical result} \\
\bottomrule
\end{tabular}
\end{table}

\subsection{Interpretation of Canonical Results}

The canonical test results show a clear progression. The enhanced Ridge model
already beats the baseline, which means the additional public-price factors are
useful. XGBoost improves further, showing that nonlinear interactions help.
The MLP family improves more than XGBoost, suggesting that standardized
continuous factors and smooth nonlinear interactions are suitable for this
cross-sectional ranking task.

The tuned MLP only slightly improves over the base MLP, so architecture tuning
alone was not the main source of improvement. The largest improvement comes
from the recent-window MLP, which changes the training distribution by using
only recent history. This model achieves \(+2.526\%\) mean 3-day excess return
on the canonical test set, exceeding the provided baseline by \(1.759\)
percentage points.

Not every attempt worked. The style-dynamic MLP does not beat the baseline on
the canonical split, so I did not select it despite later robust-split
strength. The LSTM beats the baseline only slightly, which suggests that the
available panel length and noisy 3-day label do not justify the extra sequence
model complexity.

\section{Robust Held-Out Test Results}

For the strongest and most informative candidates, I also ran the robust split
from Table \ref{tab:robust-split}. These results are shown in Table
\ref{tab:robust-results}.

\begin{table}[h]
\centering
\caption{Robust held-out test performance for selected candidates.}
\label{tab:robust-results}
\begin{tabular}{lrrr}
\toprule
Model attempt & Test excess & Positive excess rate & Rank IC \\
\midrule
Tuned MLP & \(+0.684\%\) & \(70.0\%\) & \(0.0610\) \\
Style-dynamic MLP & \(+1.130\%\) & \(70.0\%\) & \(0.0653\) \\
Model-switch router & \(+0.860\%\) & \(70.0\%\) & -- \\
Excess-target weighted MLP & \(-0.522\%\) & \(40.0\%\) & -- \\
\textbf{Recent-window MLP} & \(\mathbf{+0.804\%}\) & \(60.0\%\) & \(0.0562\) \\
\bottomrule
\end{tabular}
\end{table}

The robust split changes the interpretation of some branches. The style-dynamic
MLP is strong on the robust split but weak on the canonical split, so I judged
it unstable. The model-switch router is positive on both canonical and robust
tests but adds substantial complexity and computational cost. The excess-target
branch fails the robust split, so it is not safe despite beating the canonical
baseline. The recent-window MLP remains positive and also has the best
canonical result, which is why it is the final Stage 1 submission.

\section{Walk-Forward Test Results}

To reduce dependence on a single test window, I also used a 3-window
walk-forward evaluation for the final recent-window MLP and several research
branches. The most important aggregate result is:
\[
    \text{Recent-window MLP walk-forward excess}
    =
    +0.783\%.
\]
The weighted positive excess rate is \(66.7\%\), and the weighted rank IC is
\(-0.0251\). The walk-forward rank IC is weaker than the canonical rank IC, but
the portfolio-level excess return remains positive across multiple windows.

For the excess-target weighted MLP, the walk-forward mean excess return is only
\(+0.233\%\), much lower than the final recent-window MLP. This supports the
decision not to use the excess-target branch.

\section{Final Model Selection Based on Self-Test}

The final model selection is based on held-out test performance, not in-sample
fit. The recent-window MLP is selected because it satisfies both grading
criteria:
\begin{itemize}
    \item \textbf{Sound methodology:} it uses chronological train /
    validation / test splits, removes rows without observable 3-day targets,
    applies an embargo, selects hyperparameters only on validation, and reports
    test results only after model selection.
    \item \textbf{Performance above baseline:} it achieves \(+2.526\%\)
    canonical mean 3-day excess return versus \(+0.767\%\) for the provided
    baseline, a gain of \(+1.759\) percentage points.
\end{itemize}

The final refreshed Stage 1 submission after the official data update used
\texttt{mlp/day3\_stage1\_recent\_window/mlp\_model.py}. Its refreshed
validation-selected configuration was:
\[
    \text{lookback}=189,\qquad
    \text{model}=\texttt{narrow\_deep\_lighter\_reg},\qquad
    k=35,\qquad
    \text{weight method}=\texttt{softmax\_2.0}.
\]
The generated CSV contains 35 positive-weight stocks, total weight 1, no short
positions, and maximum single-name weight 10\%.

\section{Conclusion}

Across all attempted models, most branches exceeded the provided 3-day XGBoost
baseline on the canonical held-out test set, but only the recent-window MLP
combined the strongest canonical result with positive robust and walk-forward
checks. The self-test evidence supports the final Stage 1 choice:
\[
    \texttt{mlp/day3\_stage1\_recent\_window}.
\]

\end{document}
